\begin{definition}[The charge projection in the group-pair basis]
    \label{def:peps_regular_charge_detector_matrix}
    \lean{TNLean.PEPS.regularChargeDetectorMatrix,
        TNLean.PEPS.regularChargeTranslationMatrix}
    \leanok
    \uses{def:peps_regular_charge_translation_detector}
    Write $D_\chi$ for the matrix of the orthogonal projection onto
    $\mathcal H_\chi$ in the canonical basis of $\ell^2(G\times G)$.
    Write $T_{z,w}$ for the matrix of independent left translations.
    These are the regular-coordinate operators used in
    \cite[charge detection]{Schuch2010PEPS}, lines 2470--2486.
\end{definition}

\begin{theorem}[Matrix properties and translation covariance of the charge projection]
    \label{thm:peps_regular_charge_detector_matrix}
    \lean{TNLean.PEPS.regularChargeTranslationMatrix_eq_kronecker,
        TNLean.PEPS.regularChargeDetectorMatrix_properties,
        TNLean.PEPS.regularChargeDetectorMatrix_mulVec,
        TNLean.PEPS.regularChargeDetectorMatrix_commute_translation,
        TNLean.PEPS.regularChargeDetectorMatrix_translate_entries}
    \leanok
    \uses{def:peps_regular_charge_detector_matrix,
        thm:peps_regular_charge_detector}
    The matrix $D_\chi$ is Hermitian, positive semidefinite, and idempotent.
    It acts on coefficient vectors as the orthogonal projection onto
    $\mathcal H_\chi$. Moreover,
    \[
        T_{z,w}=U_z\otimes U_w,\qquad
        T_{z,w}D_\chi=D_\chi T_{z,w},
    \]
    and its entries satisfy
    \[
        D_\chi((zr,ws),(za,wb))=D_\chi((r,s),(a,b)).
    \]
\end{theorem}
\begin{proof}\leanok
    Express the previously constructed orthogonal projection and unitary
    translations in the canonical orthonormal basis. The two independent
    register permutations give the displayed tensor product. Comparing
    entries in their commutation identity gives the entry covariance.
\end{proof}

\begin{theorem}[Action on literal charge coefficient vectors]
    \label{thm:peps_regular_charge_detector_matrix_eigenvectors}
    \lean{TNLean.PEPS.regularChargeDetectorMatrix_chargeVector,
        TNLean.PEPS.regularChargeDetectorMatrix_other_chargeVector}
    \leanok
    \uses{thm:peps_regular_charge_detector_matrix,
        thm:peps_regular_charge_subspace_orthogonality}
    For every $p,x,y\in G$,
    \[
        D_\chi A_\chi(p;x,y)=A_\chi(p;x,y).
    \]
    If $\sigma,\tau$ are irreducible complex representations with distinct
    characters and $\sigma$ is unitary, then
    \[
        D_{\chi_\sigma}A_{\chi_\tau}(p;x,y)=0.
    \]
    These matrix identities are auxiliary to the original-spin detector in
    \cite[charge detection]{Schuch2010PEPS}, lines 2474--2486.
\end{theorem}
\begin{proof}\leanok
    Express the selected-charge and distinct-charge actions of the Hilbert-space
    projection in the canonical basis.
\end{proof}
